\documentclass[12pt]{article}

% ============================================================
% PACKAGES
% ============================================================

\usepackage{amsmath}      % Mathematical environments: align, equation, etc.
\usepackage{amsthm}       % Theorem, definition, proof environments
\usepackage{amsfonts}     % Mathematical fonts such as \mathbb
\usepackage{amssymb}      % Additional mathematical symbols
\usepackage{mathtools}    % Extensions to amsmath
\usepackage{geometry}     % Page margins
\usepackage{enumitem}     % Customize enumerate/itemize
\usepackage{mdframed}     % Framed boxes
\usepackage{tikz}         % Draw diagrams
\usepackage{hyperref}     % Clickable references and links

\geometry{
    letterpaper,
    margin=0.8in
}

\linespread{1.2}


% ============================================================
% NUMBER SYSTEMS
% ============================================================
% Instead of writing \mathbb{R} every time, use \RR.

\newcommand{\RR}{\mathbb{R}}   % Real numbers
\newcommand{\CC}{\mathbb{C}}   % Complex numbers
\newcommand{\NN}{\mathbb{N}}   % Natural numbers
\newcommand{\ZZ}{\mathbb{Z}}   % Integers
\newcommand{\QQ}{\mathbb{Q}}   % Rational numbers


% ============================================================
% COMMON MATHEMATICAL NOTATION
% ============================================================
% These commands are useful for frequently used notation.

\newcommand{\set}[1]{\left\{#1\right\}}
% Example:
%   \set{x \in \RR : x > 0}
% produces:
%   { x in R : x > 0 }

\newcommand{\abs}[1]{\left|#1\right|}
% Example:
%   \abs{x-y}
% produces:
%   |x-y|

\newcommand{\norm}[1]{\left\|#1\right\|}
% Example:
%   \norm{x}
% produces:
%   ||x||

\newcommand{\inner}[2]{\left\langle #1,#2\right\rangle}
% Example:
%   \inner{x}{y}
% produces:
%   <x,y>

\newcommand{\parens}[1]{\left(#1\right)}
% Example:
%   \parens{x+y}

\newcommand{\brac}[1]{\left[#1\right]}
% Example:
%   \brac{a,b}

\newcommand{\angles}[1]{\left\langle #1\right\rangle}
% Example:
%   \angles{x,y}


% ============================================================
% THEOREM / THEORY COMPONENTS
% ============================================================

% "definition" style is upright text.
% Use this for definitions, examples, and remarks.

\theoremstyle{definition}

% ------------------------------------------------------------
% DEFINITION
% ------------------------------------------------------------
% Use when introducing a formal mathematical definition.
%
% Usage:
%
% \begin{definition}[Metric]
% A metric on a set X is a function ...
% \end{definition}
%
% [Metric] is optional and gives the definition a name.

\newtheorem{definition}{Definition}[section]


% ------------------------------------------------------------
% EXAMPLE
% ------------------------------------------------------------
% Use to give a concrete example of a definition/theorem.
%
% Usage:
%
% \begin{example}
% Consider X = \RR with the usual distance.
% \end{example}

\newtheorem{example}[definition]{Example}


% ------------------------------------------------------------
% REMARK
% ------------------------------------------------------------
% Use for observations, comments, or small additional facts.
% It is NOT intended for major mathematical results.
%
% Usage:
%
% \begin{remark}
% This definition does not require X to be a subset of \RR^n.
% \end{remark}

\newtheorem{remark}[definition]{Remark}


% ============================================================
% THEOREM-LIKE COMPONENTS
% ============================================================

% "plain" style makes the mathematical statements italic.

\theoremstyle{plain}


% ------------------------------------------------------------
% THEOREM
% ------------------------------------------------------------
% Use for an important/general mathematical result.
%
% Usage:
%
% \begin{theorem}[Cauchy--Schwarz Inequality]
% For all x,y \in \RR^n,
% \[
% |\inner{x}{y}| \leq \norm{x}\norm{y}.
% \]
% \end{theorem}

\newtheorem{theorem}[definition]{Theorem}


% ------------------------------------------------------------
% PROPOSITION
% ------------------------------------------------------------
% Use for a useful mathematical result that is usually less
% central than a theorem.
%
% Usage:
%
% \begin{proposition}
% Every finite-dimensional subspace of \RR^n is closed.
% \end{proposition}

\newtheorem{proposition}[definition]{Proposition}


% ------------------------------------------------------------
% LEMMA
% ------------------------------------------------------------
% Use for an auxiliary result whose main purpose is to help
% prove another theorem/proposition.
%
% Usage:
%
% \begin{lemma}
% ...
% \end{lemma}

\newtheorem{lemma}[definition]{Lemma}




% ------------------------------------------------------------
% COROLLARY
% ------------------------------------------------------------
% Use for a result that follows relatively directly from a
% theorem/proposition that was just established.
%
% Usage:
%
% \begin{corollary}
% ...
% \end{corollary}

\newtheorem{corollary}[definition]{Corollary}


% ============================================================
% PROOF
% ============================================================
% The proof environment comes from amsthm.
%
% Usage:
%
% \begin{proof}
% Suppose that ...
%
% Therefore, ...
% \end{proof}
%
% LaTeX automatically adds the QED symbol at the end.


% ============================================================
% CUSTOM BOX: IDEA
% ============================================================
% Use this when you want to record intuition or the main idea
% behind a mathematical concept.
%
% Usage:
%
% \begin{idea}
% \textbf{Key idea.}
% A metric gives us a notion of distance...
% \end{idea}

\newmdenv[
    linewidth=1pt,
    skipabove=10pt,
    skipbelow=10pt
]{idea}


% ============================================================
% CUSTOM BOX: WARNING
% ============================================================
% Use this for common mistakes, distinctions, or things that
% are easy to confuse.
%
% Usage:
%
% \begin{warning}
% Do not confuse continuity with uniform continuity.
% \end{warning}

\newmdenv[
    linewidth=1pt,
    skipabove=10pt,
    skipbelow=10pt
]{warning}

% ============================================================
% CUSTOM BOX: NOTE
% ============================================================
% Use this for a short personal note, useful observation,
% connection between concepts, or explanation.
%
% Unlike "Remark", this is intended for informal notes.
%
% Usage:
%
% \begin{note}
% A sequence can be viewed as a function from $\NN$ to $\RR$.
% \end{note}

\newmdenv[
    linewidth=1pt,
    skipabove=10pt,
    skipbelow=10pt
]{note}


% ============================================================
% DOCUMENT
% ============================================================

\begin{document}


% ============================================================
% TITLE
% ============================================================

\begin{center}
    {\Large \textbf{Compact Sets}}\\[0.5em]
    {\large Lecture 11}
\end{center}

\tableofcontents

\newpage

\begin{quote}
Next best thing to being finite is compact.
\end{quote}
\section{Compact sets}

\begin{note}
Compact sets have a lot of common nice properties with finite sets.
\end{note}


\begin{definition}[Open covers]
A (open) cover of $E$ in $X$ is a collection $\{G_{\alpha}\}$ of open sets that the union of that set contains $E$.
\end{definition}

\begin{example}
We want to find cover sets of $E = (0, 1)$. We have $\{(-1,2)\}$ is a cover set, $\RR$ is a cover set, $\{V_n\}_{n=3}^{\infty} \text{ with } V_n = (\frac{1}{n}, 1-\frac{1}{n})$ is a cover set. In the last case, we can see that, we do not need every $V_n$ in order to cover $E$, but we still need to have infinitely many of them to cover $E$.
\end{example}

\begin{definition}[Subcover]
A subcover of a cover is a subcollection of that cover that still covers $E$.
\end{definition}

\begin{definition}[Compact sets]
A set $K$ in $X$ is called compact (in $X$) if \textbf{every} open cover of $K$ has a finite subcover.
\end{definition}

\begin{example}
Here are some examples:
\begin{itemize}
    \item $(0,1)$ is not compact in $\RR$, as we showed that, $\{V_n\}$ has no finite subcover.
    \item $\ZZ$ is not compact in $\RR$. The open cover with each integer endowed with an open ball, that open cover has no finite subcover.
    \item $[0,1]$ is compact.
    \item Finite set is compact.
\end{itemize}
\end{example}

\begin{note}
Showing a set is compact is not showing \textbf{a} cover has a finite subcover.
\end{note}

\subsection{Properties of compact sets}
\begin{definition}[Bounded sets]
Set $K$ is bounded if $\exists M\in\RR, p \in X, \text{ s.t. } K\in N_M(p)$.
\end{definition}

\begin{theorem}
$K$ is compact in $X$ $\Rightarrow$ $K$ is bounded.
\end{theorem}

\begin{note}
The idea for proving the theorem above is quite interesting, it construct a finite set of unit-1 balls to cover the set (given it is compact), then find the max distance between those center and let the bounding ball has the radius of max distance $+2$. Sound very barbaric, but somehow elegant.
\end{note}

\begin{theorem}
$K$ is compact in $X$, then $K$ is closed.
\end{theorem}

\section{Relative open sets}
\begin{proposition}
$Y \in X \text{ metric}$ then $Y$ also a metric space. We can say that, $Y$ inherited the metric of $X$. 
\end{proposition}

\begin{definition}[Relative nbhd]
A nbhd $N_r(x)$ in $X$ is $\{p:d(x,p)<r \text{ for } p \in X\}$.

A nbhd $N_r(x)$ in $Y$ is $\{p:d(x,p)<r \text{ for } p \in Y\}$.
\end{definition}

\begin{definition}[Relative open set]
A set $U$ is open in $Y$ if every point in $U$ is an interior point of $U$ in $Y$ (using $Y$ nbhd).
\end{definition}

\begin{theorem}
Given $E\subset Y \subset X$, then $E$ open in $Y$ $\iff$ $E = Y \cap G$ with some $G$ open in $X$.
\end{theorem}

\end{document}